\chapter{Ejemplos}
En este capítulo veremos algunos ejemplos de iteraciones de funciones racionales
para introducirnos en la rama de la Dinámica Compleja.

\section{Transformaciones de Möbius}
Empecemos con el ejemplo más sencillo, funciones racionales de grado 1. Este
tipo de funciones se conocen como transformaciones de Möbius.

\begin{definicion}[Transformación de Möbius]
    Una transformación de Möbius es una función racional de la forma
    \[
        R(z) = \frac{az+b}{cz+d}, \quad ad-bc \ne 0
    \] 
    donde $z$, $a$, $b$, $c$, $d$ son números complejos. Además, seguiremos los
    siguientes convenios con el punto del infinito: $R(\infty)=a/c$ y
    $R(-c/d)= \infty$.
\end{definicion}

Por ejemplo, analicemos la siguiente transformación de Möbius,
\[
    R(z) = \frac{3z+2}{2z+-1}.
\]

Esta función tiene un único punto fijo, $\zeta = 1$, y este es indiferente
porque $R'(1) = 1$. La pregunta que nos hacemos es: si iteramos en un punto
arbitrario del plano complejo, $z$, ¿cómo se comportará la secuencia $R^n(z)$?

Gracias a la simplicidad de las transformaciones de Möbius, es fácil encontrar
la expresión
\[
    R^n(z) = \frac{(2n+1) -2n}{2nz-(2n-1)} =  1 + \frac{z-1}{2nz-(2n-1)}.
\]
Por tanto, para todo $z \in \ExtC$, se cumple
\[
    R^n(z) \to \zeta,
\]
cuando $n \to \infty$.

Veamos otro ejemplo con un comportamiento diferente, pero igualmente sencillo:
\[
    R(z) = \frac{(1+i)z + (1-i)}{(1-i)z + (1+i)}.
\]

En este caso tenemos dos puntos fijos, $\zeta_1 = 1$ y $\zeta_2 = -1$, ambos
indiferentes (aunque en ambos ejemplos han salido puntos fijos indiferentes, no
siempre los puntos fijos de transformaciones de Möbius son indiferentes). Se
puede comprobar que
\begin{equation}\label{eq:mobiusEj2Conj}
    R(z) = g^{-1}Sg(z),
\end{equation}
con
\begin{align*}
    S(z) &= -iz, \\
    g(z) &= \frac{1+z}{1-z}.
\end{align*}
Analizando $S^n(z)$ podemos entender cómo se comporta $R^n(z)$ gracias a la
ecuación~\ref{eq:mobiusEj2Conj}. Ahora, $S^n(z)$ es muy sencilla, $S^4(z) = I(z)
= z$. Por tanto, para todo $z \in \ExtC$, se cumple
\[
    R^4(z)=z,
\]
es decir, todos los puntos forman parte de un 4-ciclo y nunca convergen a un
punto fijo (porque el 4-ciclo de estos está formado únicamente por ellos
mismos).

En general, las iteraciones de transformaciones de Möbius son muy sencillas y su
comportamiento se puede clasificar según tengan un único punto fijo o dos.

\begin{teorema}
Sea $R$ una transformación de Möbius, si $R$ tiene
\begin{enumerate}
    \item un punto fijo, $\zeta$, entonces para todo $z \in \ExtC$,
        $R^n(z) \to \zeta$ cuando $n~\to~\infty$;
    \item dos puntos fijos, entonces se tiene que
    \begin{enumerate}
        \item $R^n(z)$ converge a uno de los puntos fijos,
        \item $R^n(z)$ se mueve cíclicamente en un conjunto finito de puntos,
        \item o $R^n(z)$ forma un subconjunto denso de un círculo (o línea).
    \end{enumerate}
\end{enumerate}
\end{teorema}

\begin{proof}
    \cite[Ver][4-5]{Beardon1991}
\end{proof}